In this section, I will briefly discuss the novelty of this solution from a research perspective by discussing each of it components.

\begin{figure}[htbp]
    \centering
    \includesvg[
        width=0.95\linewidth
    ]{assets/research_components.svg}
    \caption{Major Components of the Framework}
    \label{fig:framework_components}
\end{figure}

\labelledsection{Computable Semantic Model}{The CSM formalizes the existing software engineering practices into a computable model so that it can be automatically reasoned about by an engine. By using CSP's to construct the CSM, it enables synthesis of an implementation that faithfully realizes the CSM's meaning. It is novel because this capability did not exist before and it is a necessary capability to automate the development and management of software systems. In building the CSM, I think it is important to recognize that this didn't invent meaning or processes that didn't already exist in software engineering, it simply made them executable.}

\labelledsection{Semantic Reasoning Engine}{The novelty of the engine is it uses the structure of the CSM to make the reasoning computable. Engineers will find the engine's reasoning to be particularly familiar because it mirrors their own process of reasoning about the design's meaning to develop, debug and maintain software systems. As a result, failure diagnosis is simply the engine reasoning about the world to identify missing meaning, leading to root cause analysis and the expansion of the designs meaning.}

\labelledsection{Semantic ToolKit}{The semantic toolkit enables a meaningful programming language that automates the construction of the CSM and it enables the automated development of intelligent systems that understand each other. This can be considered a breakthrough idea because this opens the door to a new automated method for constructing intelligent systems and it changes the nature of software development. The ability to automatically develop a system from its meaning is the direct result of establishing the meaning of the closed semantic as a CSM built from CSP's. Ultimately, this process has made software engineering itself executable and as a result, development and management have become automated.}

\labelledsection{Compressed Log Processor}{Domain specific compression is also a necessary capability to make the automation  of software systems possible because it overcomes a practical hurdle that prevents it. The novelty of CLP is that it proved that the preservation of the necessary information needed to automate systems management is possible at scale by exploiting the unambiguous nature of the CSM. So from a research perspective, the maturity of domain specific compression was a breakthrough that made automated systems management possible.}

However, neither the CSM, semantic toolkit, the engine or CLP solve this problem on their own. Through the elimination of ambiguity in the construction of a software system, they all exist as part of a coherent solution that works to automate the development and management of software systems. By transforming the development of software systems into a process of meaningful world building, this lays the foundation for many new areas of research.

